\section{Conclusiones}
    \begin{itemize}
        \item El conjunto \textit{Resultados únicos Saber 11} del periodo 2022-4
        es apto para el problema planteado, puesto que es público, con licencia
        abierta, tiene una variable objetivo continua y bien comportada (media
        250,1, desviación 51,8, asimetría 0,28) y reúne, para 532.792
        estudiantes, información del hogar, del colegio y de la geografía que
        existe antes de la prueba.

        \item La verificación de calidad fue tan importante como el análisis de
        las variables, dado que la mitad de las filas descargadas eran duplicados
        exactos y el puntaje global resultó ser una función determinista de los
        puntajes por área. Ambos hallazgos, de pasarse por alto, habrían
        producido métricas de modelado ficticias.

        \item Las tres hipótesis de trabajo reciben respaldo preliminar. El nivel
        socioeconómico del hogar, en particular la educación de los padres y el
        acceso a internet y computador, concentra las asociaciones más altas
        (H1); las características del colegio, sobre todo la jornada y la
        naturaleza, aportan brechas propias de 29 a 70 puntos y una interacción
        con el área (H2); y el departamento del colegio introduce una brecha de
        64 puntos entre Bogotá y Chocó que ninguna variable individual explica
        por sí sola (H3).

        \item Ninguna variable aislada correlaciona más de 0,35 con el puntaje y
        una regresión lineal aditiva explica el 24,2\,\% de la varianza con un
        error absoluto medio de 36,5 puntos. Ese es el punto de partida honesto
        del proyecto: el modelo de referencia de la fase 2 deberá reproducirlo
        con la partición definitiva y los modelos de la fase 3 deberán superarlo
        capturando las relaciones no lineales y las interacciones identificadas.

        \item Los valores faltantes no son aleatorios, quienes no responden el
        cuestionario obtienen 21,5 puntos menos, y varias variables ordinales no
        se comportan como tales (estrato, personas y cuartos del hogar). La
        estrategia de preparación debe conservar esa información en lugar de
        imputarla o descartarla, lo que define el plan de limpieza e ingeniería
        de variables de la siguiente actividad.
    \end{itemize}

    \vspace{10pt}

    Como trabajo futuro inmediato, la fase 2 implementará las doce decisiones de
    la Tabla~\ref{decisions}, construirá la partición definitiva en
    entrenamiento, validación y prueba, y establecerá el modelo de referencia; las
    fases 3 y 4 compararán modelos no lineales con ajuste de hiperparámetros,
    interpretarán el modelo final y comunicarán los resultados. Quedan abiertas
    la extensión a varios periodos y el uso de información agregada por colegio y
    municipio como variables de contexto.
